% Parameter, scalability, and case-study experiments included in the main paper.

\subsection{Parameter Study}\label{sec:exp-parameters}

We study two search-budget parameters. Beam width (beam) limits the number of partial plans retained at each search step, while the candidate limit (MC) bounds the number of completed plans considered. We report E2E latency only. In the beam sweep, MC is fixed at 8; in the MC sweep, beam is fixed at 3. The configured defaults are beam 3 and MC 4.

\begin{figure}[!htbp]
\centering
\begin{minipage}[t]{0.31\textwidth}\centering
\expresultplaceholder{(a) Beam width}
\end{minipage}\hfill
\begin{minipage}[t]{0.31\textwidth}\centering
\expresultplaceholder{(b) Candidate limit (MC)}
\end{minipage}\hfill
\begin{minipage}[t]{0.31\textwidth}\centering
\expresultplaceholder{(c) Cross-dataset summary}
\end{minipage}
\caption{E2E latency under different bounded-search parameter settings.}\label{fig:exp-parameters}
\end{figure}
\FloatBarrier

For beam width, T-Drive reaches its lowest recorded E2E latency at beam 3 (195/220 ms for simple/complex queries); wider beams do not yield consistent gains (199--205/232--237 ms). Porto is less sensitive, with recorded values between 1,043 and 1,060 ms for simple queries and 1,043 and 1,057 ms for complex queries. This suggests that retaining a wider frontier adds search effort without reliably improving the selected plan.

For MC, T-Drive improves as the limit grows from 4 to 8 (207/239 to 195/220 ms), then regresses at 10 (206/244 ms), indicating a useful balance between considering more alternatives and limiting planning overhead. Porto varies less monotonically: the simple workload is lowest at MC 10 (1,028 ms), while complex latency is similar at MC 4 and 10 (1,043 ms) and slightly higher at intermediate settings. Thus, the best MC value depends on dataset and workload difficulty; the recorded changes are in E2E latency only.

\subsection{Scalability}\label{sec:exp-scalability}

We measure E2E latency across T-Drive expansions ($\times1$, $\times10$, $\times50$, $\times100$) and Porto trajectory subsamples ($\times1$, $\times0.5$, $\times0.25$, $\times0.1$). Each scale uses the same query semantics and compares \KART{} with the five IR-level baselines. The measured latency and index-footprint results are \TBD.

\begin{figure}[!htbp]
\centering
\begin{minipage}[t]{0.31\textwidth}\centering
\expresultplaceholder{(a) T-Drive expansion}
\end{minipage}\hfill
\begin{minipage}[t]{0.31\textwidth}\centering
\expresultplaceholder{(b) Porto subsampling}
\end{minipage}\hfill
\begin{minipage}[t]{0.31\textwidth}\centering
\expresultplaceholder{(c) Data and index scale}
\end{minipage}
\caption{Query latency and data/index size across dataset scales.}\label{fig:exp-scalability}
\end{figure}
\FloatBarrier

The expected T-Drive trend is increasing E2E latency with larger replicas, since more rows and postings increase index scans and chunk retrieval. With fixed query windows and spatial coordinates, however, the candidate set need not grow in proportion to the replicated dataset; indexed plans should therefore scale more gently than FullScan. For Porto, smaller samples reduce the indexed population, candidate chunks, and exact-refinement work, so latency is expected to fall as the sample fraction decreases. At the smallest scales, fixed planning and scan startup costs may limit further gains. The measured trends are \TBD; these effects explain the expected dataset-specific behavior.

\subsection{Representative Query Case Study}\label{sec:exp-case-study}

\TBD

\begin{figure}[!htbp]
\centering
\expresultplaceholder{Representative query case study}
\caption{Case-study query, selected access plan, and result.}\label{fig:exp-case-study}
\end{figure}
\FloatBarrier
